\begin{abstract} 
\thispagestyle{plain}
\addcontentsline{toc}{chapter}{ABSTRACT}
\doublespacing
The transversity distribution function, $h_1^{q}(x)$, where $x$ is the longitudinal momentum fraction of the proton carried by quark $q$, encodes the proton's transverse spin structure at the leading twist. Extraction of it is difficult because of its chiral-odd nature. However, it can be coupled with a spin-dependent interference fragmentation function (FF), $H_1^{\sphericalangle, q}$, in polarized proton-proton ($p^\uparrow p$) collisions. The coupling of $h_1^{q}(x)$ and $H_1^{\sphericalangle, q}$ produces an experimentally measurable azimuthal correlation asymmetry, $A_{UT}$, between the spin of the fragmenting quark and the final state dihadron. A model-independent extraction of transversity from these measurements relies on the knowledge of dihadron FFs, namely the unpolarized dihadron FFs, $D_1^{h_1h_2/q(g)}$ for quarks, \emph{q} (gluons, \emph{g}). Extraction of these FFs requires measurements of the unpolarized dihadron cross section in $pp$ collisions, which are desperately needed. In $pp$ collisions, the unpolarized cross-section measurement provides access to the $D_1^{h_1h_2}$ for both quarks and gluons. This thesis outlines the measurements of the \dipion azimuthal correlation asymmetry in the mid-pseudorapidity ($|\eta|<1$) region using the RHIC run 2015 polarized $pp$ data and the measurement of the unpolarized \dipion cross-section measurement in the invariant mass bins in the forward and backward pseudorapidity regions with respect to the polarized beam using the RHIC run 2012 $pp$ data at $\sqrt{s}=200$ GeV.   
\end{abstract}
